\section{Reader's guide}

This blueprint has two simultaneous goals.  First, it follows the mathematical
architecture of the paper: exchangeability, a Hamming-slice reduction, a
two-level extremizer, a hypergeometric moment bound, and the explicit constant
$\GammaN$.  Second, it records the actual dependency path used by the checked
Lean theorem.  Those paths are displayed separately whenever they differ.

The recommended reading order is:
\begin{enumerate}[label=\arabic*.]
  \item the notation dictionary and Theorem~\ref{thm:main-mgf};
  \item the structural chain in Section~\ref{sec:paper-route};
  \item the arithmetic of $\GammaN$ in Section~\ref{sec:gamma};
  \item optimality and the variational characterization in
    Section~\ref{sec:optimality}.
\end{enumerate}

\section{Translation of the paper's notation}

\begin{definition}[Law-based exchangeability]
\label{def:exchangeability}
\lean{SharpSerfling.FinitePopulation.IsExchangeableInLaw}
\leanok
For every permutation $\pi$ of $\{1,\ldots,N\}$, the random vectors
$(X_{\pi(1)},\ldots,X_{\pi(N)})$ and $(X_1,\ldots,X_N)$ have the same
pushforward law.
\end{definition}

\begin{definition}[Zero padding and projection]
\label{def:centered-weight}
\lean{SharpSerfling.FinitePopulation.zeroPad, SharpSerfling.FinitePopulation.centeredWeight, SharpSerfling.FinitePopulation.sqNorm}
\leanok
For $w=(w_1,\ldots,w_n)$, let
\[
  \wt=(w_1,\ldots,w_n,0,\ldots,0)\in\R^N,
  \qquad
  P_{\one^\perp}z=z-\frac{\one^\top z}{N}\one.
\]
The Lean term \leanname{centeredWeight hn w} is precisely
$P_{\one^\perp}\wt$, and
\leanname{sqNorm (centeredWeight hn w)} is
$\|P_{\one^\perp}\wt\|_2^2$.
\end{definition}

\begin{definition}[Centered weighted MGF]
\label{def:exchangeable-mgf}
\lean{SharpSerfling.FinitePopulation.exchangeableMgf}
\leanok
Writing $\bar X_N=N^{-1}\sum_{j=1}^N X_j$, the formal MGF is
\[
  \E\exp\!\left\{\lambda\sum_{i=1}^n w_i(X_i-\bar X_N)\right\}.
\]
It is implemented as an integral of the exponential centered contrast on an
arbitrary probability space.
\end{definition}

\section{Main theorem and tail bound}

\begin{theorem}[Exact finite-population certificate]
\label{thm:exact-mgf-certificate}
\lean{SharpSerfling.FinitePopulation.weighted_exchangeable_mgf_centeredNorm_inLaw}
\leanok
For variables in an arbitrary interval $[a,b]$, the logarithmic MGF is bounded
with coefficient
\[
  \kappa_N\frac{N}{N-1}.
\]
This theorem is stronger than the explicit $\GammaN$ bound and is the main
kernel dependency used to close Theorem~\ref{thm:main-mgf}.
\end{theorem}

\begin{lemma}[Comparison of the exact and explicit coefficients]
\label{lem:coefficient-comparison}
\lean{SharpSerfling.ExchangeableHoeffding.kappa_le_one, SharpSerfling.ExchangeableHoeffding.Gamma_lower_variance, SharpSerfling.ExchangeableHoeffding.sharpCoefficient_le_Gamma}
\leanok
For $N\ge2$,
\[
  \kappa_N\frac{N}{N-1}\le \GammaN.
\]
The proof combines $\kappa_N\le1$ with $N/(N-1)\le\GammaN$.
\end{lemma}

\begin{theorem}[Moment-generating-function bound; Theorem 1]
\label{thm:main-mgf}
\lean{SharpSerfling.ExchangeableHoeffding.theorem_1_mgf, SharpSerfling.ExchangeableHoeffding.Gamma_isCoefficient}
\leanok
\uses{def:exchangeability,def:centered-weight,def:exchangeable-mgf,def:gamma,thm:exact-mgf-certificate,lem:coefficient-comparison}
Let $N\ge2$, let $X_1,\ldots,X_N\in[-1,1]$ be exchangeable, and let
$w\in\R^n$ with $n\le N$.  Then, for every $\lambda\in\R$,
\[
  \E\exp\!\left\{\lambda\sum_{i=1}^n w_i(X_i-\bar X_N)\right\}
  \le
  \exp\!\left\{
    \frac{\lambda^2}{2}\GammaN\|P_{\one^\perp}\wt\|_2^2
  \right\}.
\]
Equivalently, the formal declaration proves the logarithm of the left-hand
side is at most the exponent on the right.  The declaration
\leanname{Gamma_isCoefficient} strengthens the interval $[-1,1]$ statement to
the affine-normalized form valid over every interval $[a,b]$.
\end{theorem}

\begin{theorem}[Arbitrary-threshold upper tail]
\label{thm:upper-tail}
\lean{SharpSerfling.ExchangeableHoeffding.theorem_1_upperTail}
\leanok
\uses{def:centered-weight,def:gamma,lem:coefficient-comparison}
If $u>0$ and $P_{\one^\perp}\wt\ne0$, then
\[
  \Prob\!\left\{\sum_{i=1}^n w_i(X_i-\bar X_N)\ge u\right\}
  \le
  \exp\!\left\{-\frac{u^2}
  {2\GammaN\|P_{\one^\perp}\wt\|_2^2}\right\}.
\]
The implementation compares a sharper pre-existing tail certificate with
$\GammaN$.
\end{theorem}

\begin{corollary}[Confidence form; Corollary 1]
\label{cor:confidence}
\lean{SharpSerfling.ExchangeableHoeffding.corollary_1}
\leanok
\uses{thm:upper-tail}
For $0<\delta<1$ and $P_{\one^\perp}\wt\ne0$,
\[
\Prob\!\left\{
  \sum_{i=1}^n w_i(X_i-\bar X_N)
  \ge \|P_{\one^\perp}\wt\|_2
     \sqrt{2\GammaN\log(1/\delta)}
\right\}\le\delta.
\]
The nonzero hypothesis is necessary for this weak-tail formulation.  At zero
projected norm, the displayed event is $\{0\ge0\}$ and therefore has
probability one.  A strict-tail formulation would include the degenerate case.
\end{corollary}

\section{The paper's structural route}
\label{sec:paper-route}

This section records the numbered structural results that explain the proof in
the paper.  They are all checked Lean declarations.  The final theorem above is
nevertheless closed through the stronger certificate in
Theorem~\ref{thm:exact-mgf-certificate}; hence the nodes below are not falsely
marked as direct kernel dependencies of Theorem~\ref{thm:main-mgf}.

\begin{lemma}[Hoeffding's lemma; Lemma 1]
\label{lem:hoeffding}
\lean{SharpSerfling.ExchangeableHoeffding.lemma_1_hoeffding}
\leanok
If $Z\in[a,b]$ almost everywhere and $\E Z=0$, then
\[
  \E e^{\theta Z}\le \exp\!\left\{\frac{\theta^2(b-a)^2}{8}\right\}.
\]
The formal statement is measure-theoretic and assumes measurability explicitly.
\end{lemma}

\begin{lemma}[Hermite sign; Lemma 2]
\label{lem:hermite-sign}
\lean{SharpSerfling.ExchangeableHoeffding.lemma_2_hermite_sign}
\leanok
Let $x_1<x_2<x_3$ be zeros of a five-times continuously differentiable
function $h$, with $h^{(5)}>0$.  For
$p(t)=\prod_{i=1}^3(t-x_i)$,
\[
  \sum_{i=1}^3\frac{h'(x_i)}{p'(x_i)^2}>0.
\]
The ordered-node formulation is equivalent to the paper's unordered statement
after relabeling.
\end{lemma}

\begin{lemma}[Three-coordinate section; Lemma 3]
\label{lem:three-coordinate}
\lean{SharpSerfling.ExchangeableHoeffding.lemma_3_three_coordinate}
\leanok
\uses{lem:hermite-sign}
For $A,B\ge0$ with $A+B>0$, every global maximizer of
\[
  A(e^u+e^v+e^z)+B(e^{-u}+e^{-v}+e^{-z})
\]
subject to fixed $u+v+z$ and $u^2+v^2+z^2$ has two equal coordinates.
The formal statement takes the global-maximizer property as an explicit
hypothesis and concludes the disjunction of coordinate equalities.
\end{lemma}

\begin{definition}[Slice MGF and centered sphere]
\label{def:slice}
\lean{SharpSerfling.FinitePopulation.sliceMgf, SharpSerfling.FinitePopulation.centeredSphere, SharpSerfling.FinitePopulation.HasAtMostTwoValues}
\leanok
For a uniformly chosen $K$-subset $S_K$ and $y\in\R^N$,
\[
  \operatorname{sliceMgf}_{N,K}(y)
  =\E_{S_K}\exp\!\left\{\sum_{i\in S_K}y_i\right\}.
\]
The centered sphere is the set of $y$ satisfying
$\sum_i y_i=0$ and $\|y\|_2^2=\rho^2$.
\end{definition}

\begin{proposition}[Two-level maximizer; Proposition 2]
\label{prop:two-level}
\lean{SharpSerfling.ExchangeableHoeffding.proposition_2_two_level}
\leanok
\uses{def:slice,lem:three-coordinate}
On every nonempty centered fixed-radius sphere, the slice MGF has a global
maximizer with at most two distinct coordinate values.

If one value $\alpha$ has multiplicity $m$, the other value is $\beta$, and
$d=\alpha-\beta$, then centering and the norm constraint imply
\[
  \alpha=\frac{N-m}{N}d,
  \qquad
  \beta=-\frac{m}{N}d,
  \qquad
  d^2=\rho^2\frac{N}{m(N-m)}.
  \tag{2}
\]
\end{proposition}

\begin{definition}[Hypergeometric martingale factor]
\label{def:martingale-factor}
\lean{SharpSerfling.ExchangeableHoeffding.martingaleFactor}
\leanok
For $s=m\wedge(N-m)$, set
\[
  B_{N,m}=(N-s)^2\sum_{\ell=N-s}^{N-1}\ell^{-2}.
\]
\end{definition}

\begin{lemma}[Coefficient comparison for Lemma 4]
\label{lem:hypergeom-factor}
\lean{SharpSerfling.ExchangeableHoeffding.sampleVarianceFactor_le_martingaleFactor}
\leanok
\uses{def:martingale-factor}
For $1\le m\le N-1$,
\[
  \frac{m(N-m)}{N-1}\le B_{N,m}.
\]
\end{lemma}

\begin{lemma}[Hypergeometric MGF; Lemma 4]
\label{lem:hypergeom-mgf}
\lean{SharpSerfling.ExchangeableHoeffding.lemma_4_hypergeometric}
\leanok
\uses{def:martingale-factor,lem:hypergeom-factor}
If $H\sim\operatorname{Hypergeometric}(N,K,m)$ and $\mu=K/N$, then
\[
  \log\E e^{t(H-m\mu)}\le\frac{t^2}{8}B_{N,m}.
\]
The Lean proof derives this from a stronger universal hypergeometric bound and
then applies Lemma~\ref{lem:hypergeom-factor}.  Thus it checks the paper's
literal coefficient without duplicating a weaker argument.
\end{lemma}

The paper combines Proposition~\ref{prop:two-level}, equation (2), and
Lemma~\ref{lem:hypergeom-mgf} to obtain
\[
  \log\E_{S_K}\exp\!\left\{\sum_{i\in S_K}y_i\right\}
  \le\frac{\GammaN}{8}\|y\|_2^2,
  \qquad \sum_i y_i=0.
  \tag{1}
\]
This is the mathematical bridge from the Hamming slice to Theorem~1.

\section{The explicit inflation factor}
\label{sec:gamma}

\begin{definition}[Inverse-square tail and $\GammaN$]
\label{def:gamma}
\lean{SharpSerfling.ExchangeableHoeffding.inverseSquareTail, SharpSerfling.ExchangeableHoeffding.gammaTerm, SharpSerfling.ExchangeableHoeffding.Gamma}
\leanok
The reverse-indexed sum satisfies
\[
  \operatorname{inverseSquareTail}(N,s)
  =\sum_{j=0}^{s-1}(N-1-j)^{-2}
  =\sum_{\ell=N-s}^{N-1}\ell^{-2}.
\]
The formal definition is the literal finite maximum
\[
  \GammaN=
  \max_{1\le s\le\lfloor N/2\rfloor}
  \frac{N(N-s)}s\sum_{\ell=N-s}^{N-1}\ell^{-2}.
\]
Outside the paper's range $N\ge2$, the definition is set to zero.
\end{definition}

\begin{lemma}[Monotonicity of the maximizing terms]
\label{lem:gamma-mono}
\lean{SharpSerfling.ExchangeableHoeffding.inverseSquareTail_upper_half, SharpSerfling.ExchangeableHoeffding.inverseSquareTail_upper_for_monotonicity, SharpSerfling.ExchangeableHoeffding.gammaTerm_mono_step, SharpSerfling.ExchangeableHoeffding.gammaTerm_mono}
\leanok
\uses{def:gamma}
The terms in the maximum defining $\GammaN$ are nondecreasing for
$1\le s\le\lfloor N/2\rfloor$.  The key estimate is the half-integer
telescoping bound
\[
  \sum_{\ell=N-s}^{N-1}\frac1{\ell^2}
  \le
  \frac1{N-s-\tfrac12}-\frac1{N-\tfrac12}.
\]
\end{lemma}

\begin{theorem}[Maximizer and parity closed forms]
\label{thm:gamma-closed}
\lean{SharpSerfling.ExchangeableHoeffding.Gamma_eq_lastTerm, SharpSerfling.ExchangeableHoeffding.Gamma_closedForm_even, SharpSerfling.ExchangeableHoeffding.Gamma_closedForm_odd}
\leanok
\uses{lem:gamma-mono}
The maximum occurs at $s=\lfloor N/2\rfloor$, and
\[
\GammaN=
\begin{cases}
N\displaystyle\sum_{\ell=N/2}^{N-1}\ell^{-2},
&N\text{ even},\\[1em]
\dfrac{N(N+1)}{N-1}
\displaystyle\sum_{\ell=(N+1)/2}^{N-1}\ell^{-2},
&N\text{ odd}.
\end{cases}
\]
\end{theorem}

\begin{theorem}[Asymptotic expansion]
\label{thm:gamma-asymptotic}
\lean{SharpSerfling.ExchangeableHoeffding.Gamma_even_expansion_bound, SharpSerfling.ExchangeableHoeffding.Gamma_odd_expansion_bound, SharpSerfling.ExchangeableHoeffding.Gamma_asymptotic}
\leanok
\uses{thm:gamma-closed}
As $N\to\infty$,
\[
  \GammaN=1+\frac{3}{2N}+O(N^{-2}).
\]
The implementation first proves explicit remainder bounds on the even and odd
subsequences and then combines them into a filter-based big-$O$ statement.
\end{theorem}

\begin{lemma}[Comparison with the harmonic coefficient; Lemma 5]
\label{lem:barber}
\lean{SharpSerfling.ExchangeableHoeffding.harmonicNumber, SharpSerfling.ExchangeableHoeffding.barberEpsilon, SharpSerfling.ExchangeableHoeffding.lemma_5, SharpSerfling.ExchangeableHoeffding.lemma_5_eq_two}
\leanok
\uses{thm:gamma-closed}
Let
\[
  H_N=\sum_{j=1}^N\frac1j,
  \qquad
  \epsN=\frac{H_N-1}{N-H_N}.
\]
Then $\GammaN<1+\epsN$ for every $N\ge3$, while
$\Gamma_2=1+\epsilon_2$.
\end{lemma}

\section{Rate optimality and the exact variational constant}
\label{sec:optimality}

\begin{definition}[Valid exchangeable coefficient]
\label{def:coefficient}
\lean{SharpSerfling.ExchangeableHoeffding.ExchangeableHoeffdingCoefficient}
\leanok
A real number $C$ is valid at population size $N$ if the affine-normalized MGF
bound holds uniformly over every probability space, interval $[a,b]$,
exchangeable law, partial weight vector, and exponential tilt.
\end{definition}

\begin{theorem}[Exact coefficient]
\label{thm:optimal-exact}
\lean{SharpSerfling.ExchangeableHoeffding.optimalCoefficient_exact}
\leanok
\uses{def:coefficient}
The exact smallest valid coefficient is
\[
  \kappa_N\frac{N}{N-1}.
\]
The formal result proves both validity and minimality.
\end{theorem}

\begin{proposition}[Variance lower bound; Proposition 1]
\label{prop:variance-lower}
\lean{SharpSerfling.ExchangeableHoeffding.varianceLowerConstant, SharpSerfling.ExchangeableHoeffding.varianceLowerConstant_le_exact, SharpSerfling.ExchangeableHoeffding.proposition_1}
\leanok
\uses{thm:optimal-exact}
Every valid coefficient $C$ satisfies
\[
  C\ge
  \begin{cases}
    \dfrac{N}{N-1},&N\text{ even},\\[.7em]
    \dfrac{N+1}{N},&N\text{ odd}.
  \end{cases}
\]
This certifies that an inflation of order $N^{-1}$ cannot be removed.
\end{proposition}

\begin{definition}[Paper-normalized hypergeometric supremum]
\label{def:variational}
\lean{SharpSerfling.ExchangeableHoeffding.paperNormalizedLogMgf, SharpSerfling.ExchangeableHoeffding.paperVariationalValues, SharpSerfling.ExchangeableHoeffding.variationalConstant}
\leanok
For $H_{N,k,m}\sim\operatorname{Hypergeometric}(N,k,m)$ define
\[
  \psi_{N,k,m}(t)
  =\log\E\exp\!\left\{t\left(H_{N,k,m}-\frac{km}{N}\right)\right\}
\]
and
\[
  \VN=
  \max_{\substack{1\le k\le N-1\\1\le m\le N-1}}
  \frac{8N}{m(N-m)}
  \sup_{t\ne0}\frac{\psi_{N,k,m}(t)}{t^2}.
\]
The Lean definition uses the supremum of the literal set of normalized ratios.
\end{definition}

\begin{theorem}[Variational characterization; Remark 1]
\label{thm:variational}
\lean{SharpSerfling.ExchangeableHoeffding.paperNormalizedLogMgf_eq, SharpSerfling.ExchangeableHoeffding.variationalConstant_eq_exact_even, SharpSerfling.ExchangeableHoeffding.variationalConstant_eq_exact_odd, SharpSerfling.ExchangeableHoeffding.remark_1_variational}
\leanok
\uses{def:variational,thm:optimal-exact}
The variational constant equals the exact optimal exchangeable coefficient:
\[
  \VN=\kappa_N\frac{N}{N-1}=\Cstar.
\]
For odd $N$, the supremum is witnessed by an explicit nonzero tilt.  For even
$N$, the sharp value is obtained as the tilt tends to zero.
\end{theorem}

\section{Verification boundary}

All declarations attached with \leanname{\textbackslash leanok} above are
checked by the kernel.  The audit entry point prints the axioms of every public
paper result.  A separate source scan rejects proof placeholders,
project-defined axioms, and unsafe implementation escape hatches.  The expected
axiom report contains only the standard logical foundations inherited from
Lean and mathlib: \leanname{propext}, \leanname{Classical.choice}, and
\leanname{Quot.sound}.

The local paper used for this correspondence is pinned by SHA-256 digest
\begin{center}
\leanname{a7d8d51a32c619f117eb127a9a3c30b616234c0a37a4cc4c3ba44a15795a1e78}.
\end{center}
